% Week 4: From Words to Numbers — Bag of Words, TF-IDF, and Visualization
% BA2 Digital Korea - Leiden University
% Dr. Steven Denney

\documentclass[aspectratio=169,11pt]{beamer}

% ============================================
% Packages
% ============================================
\usepackage{tikz}
\usetikzlibrary{arrows.meta, positioning}
\usepackage{graphicx}
\usepackage{hyperref}
\usepackage{fontspec}
\usepackage{booktabs}
\usepackage{array}
\usepackage{tabularx}
\usepackage{amsmath}

% Theme
\usetheme{metropolis}

% Column type for wrapping text in tables
\newcolumntype{L}[1]{>{\raggedright\arraybackslash}p{#1}}

% ============================================
% Metadata
% ============================================
\title{BA2: Digital Korea}
\subtitle{Week 4: From Words to Numbers}
\author{Steven Denney}
\institute{Korean Studies \\ Leiden University}
\date{February 23, 2026}

% ============================================
% Document
% ============================================
\begin{document}

% --------------------------------------------
% Title Slide
% --------------------------------------------
\begin{frame}[plain,noframenumbering]
    \titlepage
\end{frame}

% --------------------------------------------
% Agenda
% --------------------------------------------
\begin{frame}{Today's Agenda}
    \begin{enumerate}
        \item Preprocessing review \& check-in
        \item The bag-of-words model
        \item Term frequency \& document frequency
        \item The document-term matrix
        \item TF-IDF: why we weight
        \item Orange demo: from corpus to visualization
        \item Pair activity
        \item Synthesis \& assignment
    \end{enumerate}
\end{frame}

% ============================================
\section{Review \& Check-In}
% ============================================

\begin{frame}{DataCamp: Wrangling Text}
    \framesubtitle{Introduction to Text Analysis in R --- Chapter 1}

    \textbf{Key skills you practiced:}
    \begin{itemize}
        \item \texttt{unnest\_tokens()} --- tokenize text into one-word-per-row
        \item \texttt{anti\_join(stop\_words)} --- remove common words
        \item \texttt{count(word, sort = TRUE)} --- find the most frequent words
        \item \texttt{filter()} and \texttt{group\_by()} --- compare subsets
    \end{itemize}

    \vspace{0.3cm}

    \begin{exampleblock}{Same workflow, our corpus}
        On the Data page: \texttt{week04\_text\_wrangling.R} applies these exact skills to the presidential speeches. Try it after class.
    \end{exampleblock}

    \vspace{0.3cm}

    \begin{alertblock}{Connection to today}
        DataCamp showed you that raw word counts are dominated by stopwords. Even after removing them, the most frequent words (국민, 경제) appear for \textit{every} president. Today we learn how \textbf{TF-IDF} solves this problem.
    \end{alertblock}
\end{frame}

\begin{frame}{Week 3 Check-In}
    \textbf{How did the preprocessing assignment go?} Any issues?

    \vspace{0.5cm}

    Quick review --- our preprocessing pipeline:

    \vspace{0.3cm}

    \begin{center}
    \begin{tikzpicture}[
        node distance=0.4cm,
        box/.style={draw, rounded corners, minimum width=5cm, minimum height=0.6cm, font=\small, align=center},
        arr/.style={-{Stealth}, thick}
    ]
        \node[box] (raw) {1. Raw text};
        \node[box, below=of raw] (clean) {2. Text cleaning};
        \node[box, below=of clean] (tok) {3. Tokenize + morphological analysis (Kiwi)};
        \node[box, below=of tok] (pos) {4. POS filtering (nouns, verbs, etc.)};
        \node[box, below=of pos] (stop) {5. Stopword removal};
        \node[box, below=of stop, fill=mDarkTeal!15] (result) {6. Clean content words};

        \draw[arr] (raw) -- (clean);
        \draw[arr] (clean) -- (tok);
        \draw[arr] (tok) -- (pos);
        \draw[arr] (pos) -- (stop);
        \draw[arr] (stop) -- (result);
    \end{tikzpicture}
    \end{center}
\end{frame}

\begin{frame}{Today's Question}
    \begin{center}
        \vspace{1cm}
        {\Large By end of the pipeline we have clean content words.}

        \vspace{1cm}

        {\LARGE \textbf{How do we turn words into numbers\\a computer can work with?}}

        \vspace{1cm}
    \end{center}
\end{frame}

% ============================================
\section{Bag of Words}
% ============================================

\begin{frame}{The Bag-of-Words Model}
    \textbf{The metaphor:} Cut every word from a speech, drop them into a bag, shake it up.

    \vspace{0.3cm}

    \begin{columns}[T]
        \begin{column}{0.48\textwidth}
            \textbf{What you keep}
            \begin{itemize}
                \item Word identity
                \item Word count
            \end{itemize}
        \end{column}
        \begin{column}{0.48\textwidth}
            \textbf{What you lose}
            \begin{itemize}
                \item Word order
                \item Grammar
                \item Context
            \end{itemize}
        \end{column}
    \end{columns}

    \vspace{0.4cm}

    \begin{exampleblock}{Example}
        ``국민이 경제를 걱정합니다'' vs.\ ``경제가 국민을 걱정합니다''

        \vspace{0.15cm}
        After preprocessing, both become \{국민, 경제\} --- identical in BoW.
    \end{exampleblock}
\end{frame}

\begin{frame}{Why Does BoW Work?}
    Losing word order sounds like a big sacrifice. Why is it still useful?

    \vspace{0.4cm}

    \begin{itemize}
        \item At scale (749 speeches), \textbf{statistical frequency patterns} are powerful enough to distinguish presidents, time periods, speech types
        \item The patterns we care about --- which words appear and how often --- survive the ``shake''
        \item Most computational text analysis methods build on this representation
    \end{itemize}

    \vspace{0.4cm}

    \begin{alertblock}{Textbook connection}
        Grimmer, Roberts \& Stewart Ch.\ 5 (read last week) introduced this idea.
        Ch.\ 6--7 (this week's reading) build directly on it.
    \end{alertblock}
\end{frame}

% ============================================
\section{Term Frequency \& Document Frequency}
% ============================================

\begin{frame}{Term Frequency (TF)}
    \textbf{Definition:} The count of term $t$ in document $d$.

    \vspace{0.3cm}

    \begin{exampleblock}{Activity (3 min)}
        Look at these two preprocessed excerpts (nouns only). Count the occurrences of 경제 and 국민.
    \end{exampleblock}

    \vspace{0.3cm}

    \textbf{The normalization problem:}

    Longer documents naturally have higher raw counts.

    \vspace{0.2cm}

    \textbf{Normalized TF} $= \dfrac{\text{count of term } t}{\text{total terms in document } d}$

    \vspace{0.3cm}

    \begin{columns}[T]
        \begin{column}{0.45\textwidth}
            Document A: 경제 appears 4 times in 25 tokens

            Normalized TF $= 4/25 = 0.16$
        \end{column}
        \begin{column}{0.45\textwidth}
            Document B: 경제 appears 6 times in 50 tokens

            Normalized TF $= 6/50 = 0.12$
        \end{column}
    \end{columns}
\end{frame}

\begin{frame}{Document Frequency (DF)}
    \textbf{Definition:} In how many documents does term $t$ appear?

    \vspace{0.3cm}

    \begin{columns}[T]
        \begin{column}{0.48\textwidth}
            \textbf{TF} = one word in \textbf{one} document

            \vspace{0.2cm}
            ``How often does 경제 appear in \textit{this} speech?''
        \end{column}
        \begin{column}{0.48\textwidth}
            \textbf{DF} = one word across the \textbf{corpus}

            \vspace{0.2cm}
            ``In how many speeches does 경제 appear?''
        \end{column}
    \end{columns}

    \vspace{0.5cm}

    \begin{center}
    \begin{tabular}{@{} l c l @{}}
        \toprule
        \textbf{Term} & \textbf{DF} & \textbf{Interpretation} \\
        \midrule
        국민 & $\sim$700 / 749 & Appears in almost every speech \\
        올림픽 & $\sim$15 / 749 & Appears in very few speeches \\
        \bottomrule
    \end{tabular}
    \end{center}

    \vspace{0.3cm}

    \begin{alertblock}{Question}
        Which word tells you more about what makes a speech \textit{distinctive}?
    \end{alertblock}
\end{frame}

% ============================================
\section{The Document-Term Matrix}
% ============================================

\begin{frame}{The Document-Term Matrix (DTM)}
    Rows = documents, columns = terms, cells = counts.

    \vspace{0.3cm}

    \begin{center}
    \begin{tabular}{@{} l c c c c c c @{}}
        \toprule
         & 국민 & 경제 & 민주 & 올림픽 & 통일 & 안보 \\
        \midrule
        노태우 연설 1 & 5 & 3 & 2 & 1 & 4 & 0 \\
        김영삼 연설 1 & 4 & 6 & 0 & 0 & 2 & 1 \\
        김대중 연설 1 & 3 & 4 & 5 & 0 & 1 & 0 \\
        문재인 연설 1 & 6 & 2 & 3 & 0 & 0 & 2 \\
        \bottomrule
    \end{tabular}
    \end{center}

    \vspace{0.3cm}

    \textbf{Sparsity:} Most cells are 0.

    \vspace{0.2cm}

    Our real corpus: 749 rows $\times$ thousands of columns, 95--99\% zeros. This is normal and expected.

    \vspace{0.3cm}

    \begin{block}{Key insight}
        This matrix \textbf{is} the bag-of-words representation. The computer no longer sees Korean --- it sees vectors of numbers.
    \end{block}
\end{frame}

% ============================================
\section{TF-IDF}
% ============================================

\begin{frame}{The Problem with Raw Counts}
    국민 has high TF in almost every speech.

    \vspace{0.3cm}

    \textbf{High count $\neq$ distinctive.}

    \vspace{0.3cm}

    We need a way to \textbf{down-weight} words that appear everywhere and \textbf{up-weight} words that are rare across the corpus but frequent in specific documents.

    \vspace{0.5cm}

    \begin{center}
        {\Large $\Rightarrow$ \textbf{TF-IDF}}
    \end{center}
\end{frame}

\begin{frame}{Inverse Document Frequency (IDF)}
    \textbf{IDF formula:}

    \vspace{0.2cm}

    $$\text{IDF}(t) = \log\!\left(\frac{N}{\text{DF}(t)}\right)$$

    \vspace{0.2cm}

    where $N$ = total number of documents, $\text{DF}(t)$ = number of documents containing term $t$.

    \vspace{0.4cm}

    \begin{center}
    \begin{tabular}{@{} l c c c @{}}
        \toprule
        \textbf{Term} & \textbf{DF} & \textbf{Calculation} & \textbf{IDF} \\
        \midrule
        국민 & 700 / 749 & $\log(749/700)$ & 0.03 (very low) \\
        올림픽 & 15 / 749 & $\log(749/15)$ & 1.70 (high) \\
        \bottomrule
    \end{tabular}
    \end{center}

    \vspace{0.3cm}

    Words that appear in \textit{every} document get an IDF close to 0.

    Words that appear in \textit{few} documents get a high IDF.
\end{frame}

\begin{frame}{TF-IDF: Putting It Together}
    $$\text{TF-IDF}(t, d) = \text{TF}(t, d) \times \text{IDF}(t)$$

    \vspace{0.3cm}

    Words that score highest: \textbf{frequent in one document but rare across the corpus}.

    \vspace{0.4cm}

    \begin{exampleblock}{Worked example}
        Suppose 올림픽 appears 3 times in a speech (TF = 3), and IDF = 1.70:
        \begin{itemize}
            \item TF-IDF = $3 \times 1.70 = 5.10$
        \end{itemize}

        \vspace{0.15cm}

        Meanwhile 국민 appears 8 times (TF = 8), but IDF = 0.03:
        \begin{itemize}
            \item TF-IDF = $8 \times 0.03 = 0.24$
        \end{itemize}
    \end{exampleblock}

    \vspace{0.3cm}

    \begin{block}{Write this down}
        TF-IDF answers: \textbf{what words make THIS document different from the others?}
    \end{block}
\end{frame}

\begin{frame}{Practical Note}
    \begin{itemize}
        \item Orange calculates TF-IDF \textbf{automatically} in the Bag of Words widget
        \item You do not need to compute it by hand
        \item But understanding the logic matters for \textbf{interpreting your results}
    \end{itemize}

    \vspace{0.5cm}

    \begin{columns}[T]
        \begin{column}{0.48\textwidth}
            \textbf{When a word has high TF-IDF:}
            \begin{itemize}
                \item It's frequent in this document
                \item It's rare in the corpus
                \item It's \textit{distinctive}
            \end{itemize}
        \end{column}
        \begin{column}{0.48\textwidth}
            \textbf{When a word has low TF-IDF:}
            \begin{itemize}
                \item It may be frequent, but\ldots
                \item It appears in many documents
                \item It's \textit{generic}
            \end{itemize}
        \end{column}
    \end{columns}
\end{frame}

% ============================================
% Break
% ============================================

\begin{frame}[plain]
    \begin{center}
        \vspace{2cm}
        {\Huge Break}\\[1cm]
        {\large We will resume in 5 minutes.}
        \vspace{2cm}
    \end{center}
\end{frame}

% ============================================
\section{Orange Demo}
% ============================================

\begin{frame}{Orange Workflow: From Corpus to Visualization}
    \begin{center}
    \begin{tikzpicture}[
        node distance=0.5cm and 0.8cm,
        box/.style={draw, rounded corners, minimum width=3.2cm, minimum height=0.6cm, font=\small, align=center},
        widget/.style={draw, rounded corners, minimum width=2.5cm, minimum height=0.6cm, font=\small, align=center, fill=mDarkTeal!15},
        arr/.style={-{Stealth}, thick}
    ]
        \node[box] (corpus) {Corpus};
        \node[box, below=of corpus] (pyscript) {Python Script};
        \node[box, below=of pyscript] (reload) {Corpus (reload)};
        \node[box, below=of reload] (preprocess) {Preprocess Text};
        \node[box, below=of preprocess] (bow) {Bag of Words};

        \node[widget, below left=0.6cm and -0.3cm of bow] (wc) {Word Cloud};
        \node[widget, left=0.3cm of wc] (stats) {Statistics};
        \node[widget, below right=0.6cm and -0.3cm of bow] (bar) {Bar Plot};
        \node[widget, right=0.3cm of bar] (dist) {Distributions};

        \draw[arr] (corpus) -- (pyscript);
        \draw[arr] (pyscript) -- (reload);
        \draw[arr] (reload) -- (preprocess);
        \draw[arr] (preprocess) -- (bow);
        \draw[arr] (bow) -- (wc);
        \draw[arr] (bow) -- (stats);
        \draw[arr] (bow) -- (bar);
        \draw[arr] (bow) -- (dist);
    \end{tikzpicture}
    \end{center}
\end{frame}

\begin{frame}{Demo Step 1--2: Load \& Preprocess}
    \textbf{Step 1: Corpus + Python Script} (review from Week 3)
    \begin{itemize}
        \item Load presidential speeches CSV
        \item Run Kiwi preprocessing script (morphological analysis + POS filtering)
    \end{itemize}

    \vspace{0.4cm}

    \textbf{Step 2: Preprocess Text}
    \begin{itemize}
        \item Whitespace tokenization (our script already split morphemes with spaces)
        \item Load custom stopword list
    \end{itemize}
\end{frame}

\begin{frame}{Demo Step 3: Bag of Words Widget}
    The Bag of Words widget converts text to a document-term matrix.

    \vspace{0.3cm}

    \textbf{Weighting options:}
    \begin{center}
    \begin{tabular}{@{} l l @{}}
        \toprule
        \textbf{Option} & \textbf{What it does} \\
        \midrule
        Count & Raw term frequency \\
        Binary & 1 if present, 0 if absent \\
        Sublinear & $1 + \log(\text{TF})$ --- dampens high counts \\
        TF-IDF & Term frequency $\times$ inverse document frequency \\
        \bottomrule
    \end{tabular}
    \end{center}

    \vspace{0.3cm}

    We'll select \textbf{TF-IDF}.

    \vspace{0.2cm}

    Open the data table output: you'll see the actual DTM.

    749 rows $\times$ thousands of columns. Most cells are 0 --- this is the sparsity we discussed.
\end{frame}

\begin{frame}{Demo Step 4: Visualization Widgets}
    \begin{columns}[T]
        \begin{column}{0.48\textwidth}
            \textbf{Word Cloud}
            \begin{itemize}
                \item Visual overview of word frequencies
                \item Filter by president to compare
                \item Bigger word = higher weight
            \end{itemize}

            \vspace{0.4cm}

            \textbf{Statistics}
            \begin{itemize}
                \item Sort by TF-IDF scores
                \item Compare: words that rank high by TF but low by TF-IDF, and vice versa
            \end{itemize}
        \end{column}
        \begin{column}{0.48\textwidth}
            \textbf{Bar Plot}
            \begin{itemize}
                \item Top 20 terms by TF-IDF for a selected president
                \item Compare across presidents
            \end{itemize}

            \vspace{0.4cm}

            \textbf{Distributions}
            \begin{itemize}
                \item Select a term (e.g., 경제)
                \item See distribution across presidents
                \item Compare: 경제 vs.\ 민주
            \end{itemize}
        \end{column}
    \end{columns}
\end{frame}

% ============================================
\section{Pair Activity}
% ============================================

\begin{frame}{Guided Pair Activity (10 min)}
    Work in pairs with the shared \texttt{.ows} workflow file.

    \vspace{0.4cm}

    \textbf{Task (5 min):}
    \begin{enumerate}
        \item Switch Bag of Words between \textbf{Count} and \textbf{TF-IDF}
        \item Compare the Word Clouds. Write down:
        \begin{itemize}
            \item One word that's prominent with Count but fades with TF-IDF
            \item One word that becomes \textit{more} prominent with TF-IDF
        \end{itemize}
        \item Open Distributions for 경제 --- which president uses it most?
    \end{enumerate}

    \vspace{0.4cm}

    \textbf{Debrief (5 min):} 2--3 pairs share findings.

    \vspace{0.2cm}

    \begin{alertblock}{Connect back}
        Why did the word clouds change? Explain using the IDF formula.
    \end{alertblock}
\end{frame}

% ============================================
\section{Synthesis \& Assignment}
% ============================================

\begin{frame}{Today's Arc}
    \begin{center}
    \begin{tikzpicture}[
        node distance=0.35cm,
        box/.style={draw, rounded corners, minimum width=2.8cm, minimum height=0.55cm, font=\small, align=center},
        arr/.style={-{Stealth}, thick}
    ]
        \node[box] (pre) {Preprocessing};
        \node[box, below=of pre] (bow) {Bag of Words};
        \node[box, below=of bow] (tf) {Term Frequency};
        \node[box, below=of tf] (df) {Document Frequency};
        \node[box, below=of df] (dtm) {Document-Term Matrix};
        \node[box, below=of dtm] (tfidf) {TF-IDF};
        \node[box, below=of tfidf, fill=mDarkTeal!15] (viz) {Visualization};

        \draw[arr] (pre) -- (bow);
        \draw[arr] (bow) -- (tf);
        \draw[arr] (tf) -- (df);
        \draw[arr] (df) -- (dtm);
        \draw[arr] (dtm) -- (tfidf);
        \draw[arr] (tfidf) -- (viz);
    \end{tikzpicture}
    \end{center}
\end{frame}

\begin{frame}{Assignment}
    \textbf{Reproduce the demo workflow:}

    \vspace{0.2cm}

    \begin{enumerate}
        \item Corpus $\rightarrow$ Python Script $\rightarrow$ Preprocess Text $\rightarrow$ Bag of Words (TF-IDF)
        \item Connect at least \textbf{two visualization widgets}
        \item Take a screenshot of your workflow
    \end{enumerate}

    \vspace{0.3cm}

    \textbf{Submit:}
    \begin{itemize}
        \item Create a \texttt{week04/} folder in your repository
        \item Add your \texttt{.ows} file + screenshot
        \item Commit and push via GitHub Desktop
    \end{itemize}

    \vspace{0.4cm}

    \begin{block}{Readings for this week}
        Grimmer, Roberts \& Stewart --- Chapters 6--7.

        Optional: Bollen et al.\ (2021) on cognitive distortions in language.
    \end{block}
\end{frame}

\begin{frame}{Looking Ahead}
    \textbf{Week 5:} Practice \& Deepen --- Hands-On Lab
    \begin{itemize}
        \item Extended practice with BoW/TF-IDF workflows
        \item Comparing presidents, interpreting visualizations
        \item Come with your assignment done
    \end{itemize}

    \vspace{0.3cm}

    \textbf{Week 6:} Midterm Review \& Assessment
    \begin{itemize}
        \item Everything from Weeks 1--5
    \end{itemize}

    \vspace{0.5cm}

    \begin{center}
        {\Large Thank you!}\\[0.3cm]
        \normalsize\texttt{s.c.denney@hum.leidenuniv.nl}
    \end{center}
\end{frame}

\end{document}
